\documentclass[11pt]{amsart}
\usepackage[margin=1in]{geometry}
\usepackage{mathptmx}
\usepackage{amssymb}
\usepackage{xcolor}
\usepackage{enumitem}
\usepackage{xurl}
\usepackage[colorlinks=true,linkcolor=blue!50!black,citecolor=blue!50!black,urlcolor=blue!50!black]{hyperref}
\usepackage{microtype}

\newtheorem{theorem}{Theorem}[section]
\newtheorem{lemma}[theorem]{Lemma}
\newtheorem{corollary}[theorem]{Corollary}
\newtheorem{proposition}[theorem]{Proposition}
\theoremstyle{definition}
\newtheorem{conjecture}[theorem]{Conjecture}
\theoremstyle{remark}
\newtheorem{remark}[theorem]{Remark}

\newcommand{\cL}{\mathcal{L}}
\newcommand{\cS}{\mathcal{S}}
\newcommand{\cV}{\mathcal{V}}
\newcommand{\cO}{\mathcal{O}}
\newcommand{\Q}{\mathbb{Q}}
\newcommand{\Z}{\mathbb{Z}}
\newcommand{\WP}{\mathrm{WP}}

% The archive id of the strings paper on hexagonmath.org and the release of this note.
\newcommand{\hexagonid}{\href{https://hexagonmath.org/2610.00446v1}{2610.00446v1}}
\newcommand{\ownrelease}{https://github.com/creelie/geometric-selection-conjecture/releases/tag/v1.0.0}

\title{The Mirror Clause of the Geometric Selection Conjecture Is Empty}

\author{Deep Bhattacharjee}
\address{Formerly, Electro-Gravitational Space Propulsion Laboratory (EGSPL), Bhubaneswar, Odisha 751030, India}
\address{Formerly, University of Calcutta, Senate House, 87/1 College Street, Kolkata 700073, West Bengal, India}
\email{itsdeep@live.com}
\urladdr{\href{https://orcid.org/0000-0003-0466-750X}{https://orcid.org/0000-0003-0466-750X}}
\thanks{Corresponding author; second e-mail address \nolinkurl{d.bhattacharjee@erl-forschung.de}}

\author{Ushashi Bhattacharya}
\address{Formerly, National Taiwan University, Taipei 106319, Taiwan}
\email{rai.bhattacharya.in@gmail.com}
\urladdr{\href{https://orcid.org/0009-0002-2254-3914}{https://orcid.org/0009-0002-2254-3914}}

\subjclass[2020]{Primary 81T30, 14J32; Secondary 14J33, 32G20}
\keywords{heterotic string, Calabi--Yau threefold, mirror symmetry, generations, Geometric Selection Conjecture}

\begin{document}

\begin{abstract}
The Geometric Selection Conjecture asks for a functional on the heterotic
landscape, built from the Weil--Petersson metric and invariant under mirror
symmetry, whose unique minimum is a phenomenologically viable vacuum. It was
shown in our paper on the Dwork quintic that the functional carries no
information beyond its minimum: the conjecture holds if and only if there is a
viable vacuum fixed by the mirror involution wherever that involution is
defined. We prove that the fixed-point clause is never satisfied. A vacuum fixed
by the mirror involution has a threefold isomorphic to its own mirror, hence
equal Hodge numbers $h^{1,1}=h^{2,1}$, vanishing Euler number and, on its
component of the landscape, a bundle with $c_3=0$; its net number of families is
therefore zero and it violates the three-family condition. Consequently no
vacuum built on the tangent bundle of any Calabi--Yau threefold, or on any
deformation of it, can be selected, and the conjecture is equivalent to the
bare statement that one viable vacuum exists on a component of the landscape
not of standard-embedding type, with no condition involving mirror
symmetry. The same holds when invariance under any compact group of
isometries containing the mirror involution is demanded. The existence
statement itself remains open.
\end{abstract}

\maketitle

\section{Introduction}\label{sec:intro}

The heterotic string on a product of four-dimensional space-time with a
Calabi--Yau threefold $X$ and a holomorphic bundle $V$ on $X$ gives a
four-dimensional gauge theory whose charged chiral multiplets are counted by
the cohomology of $V$~\cite{CHSW85,Wit85}. A conjecture of Bhattacharjee,
Bhattacharya and Ghosh~\cite{Bha26a}, the Geometric Selection Conjecture,
proposes that the vacuum of the real world is singled out among all such pairs
$(X,V)$ by a functional constructed from the natural metric on their moduli,
proper on every component, with a unique minimum, invariant under mirror
symmetry, and attaining its minimum at a vacuum that meets nine
phenomenological conditions (C1)--(C9). In~\cite{BB26} we proved that the
functional is determined by the vacuum it selects: the conjecture holds if and
only if there is a viable vacuum $\cV^\star$ with $\mu(\cV^\star)=\cV^\star$
whenever $\cV^\star$ lies in the domain $\cL^\mu$ of the mirror involution
$\mu$ (Theorem~D there, proved as Theorem~22). The content of the conjecture is therefore an
existence statement with one clause attached to it, the clause on mirror
symmetry.

This note removes that clause. The argument is short. A vacuum fixed by $\mu$
is isomorphic to its mirror, so its threefold has $h^{1,1}=h^{2,1}$ and Euler
number zero. The domain $\cL^\mu$ is made of the components of the landscape
that contain a vacuum of the standard embedding $V=T_X$, on which the Chern
classes of $V$ are constant, so $c_3(V)=c_3(T_X)=\chi(X)=0$ there. The index
theorem then gives zero net families, and condition (C2), three net families,
fails. So no fixed point of $\mu$ is viable, every viable vacuum lies outside
$\cL^\mu$, and the clause is vacuous.

\begin{theorem}\label{thm:main}
The Geometric Selection Conjecture holds if and only if the heterotic
landscape contains a viable vacuum outside the domain $\cL^\mu$ of the mirror
involution, that is, on a component that contains no vacuum of the standard
embedding. No vacuum of the form $(X,V)$ with $V$ a deformation of the
tangent bundle of $X$, on any Calabi--Yau threefold $X$ or any free quotient
of one, is the minimum of a functional with the properties the conjecture
demands, whether or not it is viable.
\end{theorem}

Remark~24 of~\cite{BB26} already noted that for a viable vacuum off the
standard embedding the clause on $\mu$ is empty. The new point is that the
clause can never be met inside $\cL^\mu$, so the vacua of standard type drop
out of the conjecture altogether and what remains is existence on the other
components. Section~\ref{sec:setting} recalls the setting
of~\cite{Bha26a,BB26}, Section~\ref{sec:proof} proves the theorem, and
Section~\ref{sec:remains} states what is left of the conjecture.

\section{The landscape, the mirror involution and viability}\label{sec:setting}

We use the setting of~\cite{Bha26a} as recalled in~\cite[Section~9]{BB26}.
A heterotic vacuum is a pair $(X,V)$ of a Calabi--Yau threefold $X$, that is a
compact K\"ahler threefold with $K_X\cong\cO_X$ and $h^{1,0}=h^{2,0}=0$, and a
slope-polystable holomorphic vector bundle $V$ on $X$ with $c_1(V)=0$ for which
$c_2(T_X)-c_2(V)$ is effective. The landscape $\cL$ is the set of isomorphism
classes of vacua, where $(X,V)\cong(X',V')$ when there is an isomorphism
$X\to X'$ carrying $V'$ to $V$. It is the disjoint union of its components
$\cL_\tau$ of fixed topological type $\tau$; the type records the
diffeomorphism type of $X$ together with the rank and the Chern classes of
$V$. Each component carries the metric $g=G^{\WP}\oplus G^V$ made of the
Weil--Petersson metric on the complex structures of $X$ and the $L^2$ metric
on the bundle moduli. Two facts about a component are used below: the Euler
number $\chi(X)=2\big(h^{1,1}(X)-h^{2,1}(X)\big)$ is constant on it, since $X$
only deforms, and the Chern classes of $V$ are constant on it, since they are
part of $\tau$.

The \emph{standard embedding} is the vacuum $(X,T_X)$. Its component is the
component of \emph{standard type} of $X$; it contains the deformations of
$T_X$ as a polystable bundle, and on all of it $c_3(V)=c_3(T_X)=\chi(X)$.

The mirror involution $\mu$ is defined on the union $\cL^\mu$ of the
components of standard type and sends a vacuum to its mirror: the threefold of
$\mu(X,V)$ is a mirror $\widehat X$ of $X$, so that
\[
h^{1,1}(\widehat X)=h^{2,1}(X),\qquad h^{2,1}(\widehat X)=h^{1,1}(X),
\]
and $\mu$ is an isometry of $g$. This is the involution of~\cite{Bha26a},
recorded there for the Weil--Petersson metric, and it is the mirror map of
$(2,2)$ vacua of~\cite{GP90,CDGP91} extended along the components they lie
on; nothing below depends on how the extension is made, only on the two
displayed equalities for the underlying threefolds.

A vacuum is \emph{viable} when it satisfies the nine conditions (C1)--(C9)
of~\cite{Bha26a}. Only the second is used here:
\begin{enumerate}[label=\textup{(C2)},leftmargin=*]
\item the vacuum has three net families, that is, the index of the Dirac
operator coupled to $V$ on $X$, counted with the Wilson line if $X$ is a free
quotient, equals three in absolute value.
\end{enumerate}
In every case of the heterotic dictionary the index in question is the
holomorphic Euler characteristic $\chi(X,V)=\sum_q(-1)^q h^q(X,V)$, up to
sign; for the standard embedding this is~\cite[Theorem~6]{BB26} and in general
it is the index formula of~\cite[Chapter~14]{GSW87}. A Wilson line on a free
quotient $Y=X/H$ breaks the gauge group but does not change the index of $V_Y$
on $Y$, which is $\chi(X,V)/|H|$; and $(Y,V_Y)$ is itself a point of $\cL$,
so it is enough to treat an arbitrary vacuum and let $X$ stand for $Y$.

The conjecture, in the form stated in~\cite{Bha26a} and recalled
in~\cite{BB26}, reads as follows.

\begin{conjecture}[Geometric Selection Conjecture]\label{conj:gsc}
There is a functional $\cS:\cL\to[0,\infty)$, constructed from the metric
$g$, such that \textup{(G1)} $\cS$ is smooth and proper on each component
$\cL_\tau$; \textup{(G2)} $\cS$ has a unique global minimum
$\cV^\star\in\cL$; \textup{(G3)} $\cV^\star$ is viable; \textup{(G4)}
$\cS\circ\mu=\cS$ on $\cL^\mu$.
\end{conjecture}

\begin{theorem}[{\cite[Theorem~22]{BB26}}]\label{thm:D}
Conjecture~\ref{conj:gsc} holds if and only if there is a viable vacuum
$\cV^\star$ with $\mu(\cV^\star)=\cV^\star$ whenever $\cV^\star\in\cL^\mu$.
\end{theorem}

\section{Fixed points of the mirror involution have no families}\label{sec:proof}

\begin{lemma}\label{lem:index}
Let $X$ be a Calabi--Yau threefold and $V$ a holomorphic vector bundle on $X$
with $c_1(V)=0$. Then $\chi(X,V)=\tfrac12\int_Xc_3(V)$.
\end{lemma}

\begin{proof}
By Hirzebruch--Riemann--Roch, $\chi(X,V)=\int_X\operatorname{ch}(V)\,
\operatorname{td}(X)$, and in degree six the integrand is
$\operatorname{ch}_3(V)+\operatorname{ch}_1(V)\operatorname{td}_2(X)+
\operatorname{rk}(V)\operatorname{td}_3(X)$. Here
$\operatorname{ch}_1(V)=c_1(V)=0$, $\operatorname{td}_3(X)=
\tfrac1{24}c_1(X)c_2(X)=0$ since $c_1(X)=0$, and
$\operatorname{ch}_3(V)=\tfrac16\big(c_1^3-3c_1c_2+3c_3\big)(V)=
\tfrac12c_3(V)$.
\end{proof}

\begin{lemma}\label{lem:fixed}
Let $(X,V)\in\cL^\mu$ satisfy $\mu(X,V)=(X,V)$. Then $h^{1,1}(X)=h^{2,1}(X)$,
$\chi(X)=0$, $c_3(V)=0$ and $\chi(X,V)=0$. In particular $(X,V)$ has no net
families and is not viable.
\end{lemma}

\begin{proof}
Since $\mu(X,V)=(X,V)$ in $\cL$, the two vacua are isomorphic, so the
threefold $\widehat X$ of $\mu(X,V)$ is isomorphic to $X$. Therefore
$h^{1,1}(X)=h^{1,1}(\widehat X)=h^{2,1}(X)$, and
$\chi(X)=2\big(h^{1,1}(X)-h^{2,1}(X)\big)=0$. The vacuum lies in $\cL^\mu$, so
its component is of standard type and contains a vacuum $(X_0,T_{X_0})$ with
$X_0$ a deformation of $X$. Chern classes are constant on the component, so
$c_3(V)=c_3(T_{X_0})=\chi(X_0)$, and Euler numbers are constant on it, so
$\chi(X_0)=\chi(X)=0$. Hence $c_3(V)=0$, and $\chi(X,V)=0$ by
Lemma~\ref{lem:index}. The net number of families is $|\chi(X,V)|=0\ne3$, so
(C2) fails.
\end{proof}

\begin{proof}[Proof of Theorem~\ref{thm:main}]
By Theorem~\ref{thm:D} the conjecture is equivalent to the existence of a
viable vacuum $\cV^\star$ that is fixed by $\mu$ if it lies in $\cL^\mu$. By
Lemma~\ref{lem:fixed} a vacuum in $\cL^\mu$ fixed by $\mu$ is not viable. So a
viable $\cV^\star$ satisfies the clause exactly when $\cV^\star\notin\cL^\mu$,
and the conjecture holds if and only if there is a viable vacuum outside
$\cL^\mu$. This is the first statement.

For the second, let $\cS$ be any functional satisfying (G1)--(G4) and let
$\cV^\star$ be its minimum. If $\cV^\star\in\cL^\mu$ then
$\cS(\mu\cV^\star)=\cS(\cV^\star)$ by (G4), so $\mu\cV^\star$ is also a global
minimum and $\mu\cV^\star=\cV^\star$ by (G2); by Lemma~\ref{lem:fixed}
$\cV^\star$ then has no families, against (G3). So the minimum of every such
functional lies outside $\cL^\mu$, and no vacuum of standard type, which is to
say no $(X,V)$ with $V$ a deformation of $T_X$ on $X$ or on a free quotient, is
ever selected. This holds whether or not such a vacuum is viable: a vacuum of
standard type has $|\chi(X)|/2$ net families by Lemmas~\ref{lem:index}
and~\ref{lem:fixed}, and the quotient of the Tian--Yau threefold by its free
$\Z/3$ has $\chi=-6$, so three families of standard type do occur in the
landscape; they are never the minimum.
\end{proof}

\begin{corollary}\label{cor:two}
If mirror pairs are identified in $\cL$, each pair counted as a single point
on the ground that both members define the same superconformal field theory,
then $\mu$ is the identity on $\cL^\mu$, condition \textup{(G4)} is empty, and
Conjecture~\ref{conj:gsc} holds if and only if there is a viable vacuum of any
type.
\end{corollary}

\begin{proof}
When mirror pairs are one point every point of $\cL^\mu$ is fixed by $\mu$, the
clause of Theorem~\ref{thm:D} is satisfied by every vacuum, and
Theorem~\ref{thm:D} gives the equivalence with existence. Lemma~\ref{lem:fixed}
does not apply, because a point that is a pair is a class of two vacua, not a
vacuum isomorphic to its mirror.
\end{proof}

\begin{remark}[The compact-group reading]\label{rem:group}
In~\cite[Remark~23]{BB26} the words ``constructed from the metric'' were
read as invariance under a compact group $G$ of isometries of $(\cL,g)$
containing $\mu$, and the conjecture in that reading was shown equivalent to
the existence of a viable vacuum fixed by $G$ wherever $G$ acts. A $G$-fixed
vacuum in $\cL^\mu$ is $\mu$-fixed, so by Lemma~\ref{lem:fixed} it is not
viable, and the same argument gives: in the compact-group reading the
conjecture holds if and only if there is a viable vacuum outside $\cL^\mu$
fixed by $G$ wherever $G$ acts on the components outside $\cL^\mu$. If $G$
acts only on $\cL^\mu$, as $\mu$ does, this is again the bare existence of a
viable vacuum outside $\cL^\mu$.
\end{remark}

\begin{remark}[Self-mirror threefolds]
Lemma~\ref{lem:fixed} does not say that a fixed point of $\mu$ is impossible,
only that it is not viable. Threefolds isomorphic to their mirrors exist; the
Schoen threefold with $h^{1,1}=h^{2,1}=19$ and its free quotient by
$\Z/3\times\Z/3$ with $h^{1,1}=h^{2,1}=3$ are examples, and on each of them the
standard embedding has zero net families, as the lemma requires. The vacua of
the standard embedding on the quintic pencil, on the three weighted pencils of
Fermat type and on their free quotients, which were studied in~\cite{BB26},
have $\chi\ne0$ and are therefore never fixed by $\mu$; they fail (C2) for
the independent reason computed there, namely that $|\chi|/2$ is $100/|H|$ on
quintic quotients and never three on the weighted pencils.
\end{remark}

\section{What remains of the conjecture}\label{sec:remains}

After Theorem~\ref{thm:main} the Geometric Selection Conjecture is the
sentence: \emph{there is a heterotic vacuum, on a component that contains no
standard embedding, satisfying \textup{(C1)}--\textup{(C9)}}. Nothing in it refers to a
functional, to the Weil--Petersson metric or to mirror symmetry. The
conditions divide into two groups.

Conditions (C1)--(C6), anomaly cancellation and stability, three net
families, a Wilson line breaking to the Standard Model group, three families
without chiral exotics, masses for the vectorlike exotics and doublet--triplet
splitting, are satisfied by known vacua: (C1)--(C4) by the heterotic standard
models of Braun, He, Ovrut and Pantev~\cite{BHOP05} and of Bouchard and
Donagi~\cite{BD06}, and (C1)--(C6) by line-bundle models of Anderson, Gray,
Lukas and Palti~\cite{AGLP11,AGLP12}. All of these have $V$ of rank four or
five with $c_3(V)\ne\chi(X)$, so they lie outside $\cL^\mu$, as
Theorem~\ref{thm:main} requires of any candidate.

Conditions (C7)--(C9), stabilised moduli, an acceptable rank and hierarchy of
the Yukawa masses, and consistent supersymmetry breaking, are stated
in~\cite{Bha26a} by comparison with measurement. For none of the vacua above
are all three established, and until the three conditions are given as
inequalities with named constants there is no mathematical statement to
prove or refute in them. The conjecture therefore stands exactly where
Theorem~\ref{thm:main} leaves it: equivalent to the existence of one vacuum
satisfying (C1)--(C9), with the knowledge that the vacuum, if it exists, is
not a deformation of the standard embedding, and with its remaining content
entirely in the three conditions that are not yet mathematical.

\section*{Data availability}

No data were generated. The source of this note and a program that checks the
lemma and the corollary on the Hodge numbers of self-mirror and three-family
threefolds are in release 1.0.0 of the repository
\url{https://github.com/creelie/geometric-selection-conjecture}
(\url{\ownrelease}); the same check is the module
\texttt{construct/gsc\_mirror\_fixed.py} of the Anaxia engine,
\url{https://github.com/creelie/eBLLM}.

\section*{Declarations}

\textbf{Competing interests.} The authors declare no competing interests.
\textbf{Funding.} No funding was received.
\textbf{Use of AI tools.} The authors used Claude Code (Anthropic) to draft
the text of this note and the checking module from the authors' argument;
the authors verified every statement and take full responsibility for the
content.

\emergencystretch=1.5em
\begin{thebibliography}{AGLP12}\raggedright

\bibitem{AGLP11}
L.~B.~Anderson, J.~Gray, A.~Lukas and E.~Palti, Two hundred heterotic
standard models on smooth Calabi--Yau threefolds, Phys. Rev. D \textbf{84}
(2011), 106005,
\href{https://doi.org/10.1103/PhysRevD.84.106005}{doi:10.1103/PhysRevD.84.106005}.

\bibitem{AGLP12}
L.~B.~Anderson, J.~Gray, A.~Lukas and E.~Palti, Heterotic line bundle
standard models, J. High Energy Phys. \textbf{2012}, no.~6, 113,
\href{https://doi.org/10.1007/JHEP06(2012)113}{doi:10.1007/JHEP06(2012)113}.

\bibitem{Bha26a}
D.~Bhattacharjee, U.~Bhattacharya and S.~Ghosh, One missing axiom for the
complete closure of the heterotic landscape: resonance classifier of a
vibrating string creates the typical particle classes, Preprints.org (2026),
\href{https://doi.org/10.20944/preprints202603.0792.v2}{doi:10.20944/preprints202603.0792.v2}.

\bibitem{BB26}
D.~Bhattacharjee and U.~Bhattacharya, Which string vibrations become which
generations on the Dwork quintic, hexagonmath.org (2026), archive
\hexagonid; code and checks at
\url{https://github.com/creelie/strings/releases/tag/v1.2.0}.

\bibitem{BD06}
V.~Bouchard and R.~Donagi, An $SU(5)$ heterotic standard model, Phys. Lett.
B \textbf{633} (2006), 783--791,
\href{https://doi.org/10.1016/j.physletb.2005.12.042}{doi:10.1016/j.physletb.2005.12.042}.

\bibitem{BHOP05}
V.~Braun, Y.-H.~He, B.~A.~Ovrut and T.~Pantev, A heterotic standard model,
Phys. Lett. B \textbf{618} (2005), 252--258,
\href{https://doi.org/10.1016/j.physletb.2005.05.007}{doi:10.1016/j.physletb.2005.05.007}.

\bibitem{CDGP91}
P.~Candelas, X.~C.~de la Ossa, P.~S.~Green and L.~Parkes, A pair of
Calabi--Yau manifolds as an exactly soluble superconformal theory, Nucl.
Phys. B \textbf{359} (1991), 21--74,
\href{https://doi.org/10.1016/0550-3213(91)90292-6}{doi:10.1016/0550-3213(91)90292-6}.

\bibitem{CHSW85}
P.~Candelas, G.~T.~Horowitz, A.~Strominger and E.~Witten, Vacuum
configurations for superstrings, Nucl. Phys. B \textbf{258} (1985), 46--74,
\href{https://doi.org/10.1016/0550-3213(85)90602-9}{doi:10.1016/0550-3213(85)90602-9}.

\bibitem{GP90}
B.~R.~Greene and M.~R.~Plesser, Duality in Calabi--Yau moduli space, Nucl.
Phys. B \textbf{338} (1990), 15--37,
\href{https://doi.org/10.1016/0550-3213(90)90622-K}{doi:10.1016/0550-3213(90)90622-K}.

\bibitem{GSW87}
M.~B.~Green, J.~H.~Schwarz and E.~Witten, \emph{Superstring Theory, Volume 2:
Loop Amplitudes, Anomalies and Phenomenology}, Cambridge University Press,
Cambridge, 1987.

\bibitem{Wit85}
E.~Witten, Symmetry breaking patterns in superstring models, Nucl. Phys. B
\textbf{258} (1985), 75--100,
\href{https://doi.org/10.1016/0550-3213(85)90603-0}{doi:10.1016/0550-3213(85)90603-0}.

\end{thebibliography}

\end{document}
